\chapter{Tres años, 1959--1962}
\label{ch:revolution}

\section{La cuestión del viraje}

La pregunta más discutida de la historia cubana moderna es si Fidel Castro era en
enero de 1959 un comunista que lo ocultaba, o un nacionalista radical al que la
reacción norteamericana empujó al campo soviético. Sesenta años de trabajo
académico no lo han resuelto y este libro tampoco lo hará, pero conviene enunciar
la forma de la evidencia, porque de la lección que se extraiga depende el
tratamiento que la Parte III dé a la relación con Estados Unidos.

Contra la tesis del ocultamiento: el programa publicado del Movimiento 26 de
Julio era nacionalista y constitucionalista; el partido comunista se opuso a la
insurrección casi hasta el final y los guerrilleros desconfiaban de él por eso;
el primer gabinete fue liberal, encabezado por el juez Manuel Urrutia como
presidente y el abogado José Miró Cardona como primer ministro, con figuras como
Felipe Pazos ---un economista ortodoxo--- en el banco central; y Castro pasó
abril de 1959 en Estados Unidos asegurando a sus auditorios que no era comunista,
en un momento en que decirlo no tenía ninguna ventaja táctica evidente que no
pudiera haber obtenido más barato.

A favor: Raúl Castro y el Che Guevara eran marxistas mucho antes de 1959 y
ninguno era una figura marginal; las estructuras paralelas que desplazaron al
gabinete formal estaban en pie en cuestión de meses; el ritmo al que se removió a
los ministros liberales ---Urrutia fuera en julio de 1959, Pazos en noviembre,
Huber Matos detenido en octubre por objetar la influencia comunista en el
ejército y condenado a veinte años--- era demasiado rápido para haber sido
improvisado; y el propio Castro dijo en 1961 que había sido marxista-leninista
todo el tiempo, aunque en 1961 dijo muchas cosas.

La lectura más defendible es que ninguna de las dos versiones puras sobrevive al
registro \citep{farber2006}. La dirigencia revolucionaria contenía marxistas
convencidos y liberales convencidos, el líder no era ninguna de las dos cosas y
era sobre todo un nacionalista con un instinto extraordinario para saber dónde
estaba el poder, y la secuencia de 1959--60 fue una interacción genuina en la que
cada movimiento de una parte abarataba el siguiente movimiento de la otra. La
reforma agraria amenazaba las propiedades norteamericanas; la respuesta
norteamericana amenazaba la reforma; la alternativa soviética volvía sobrevivible
esa respuesta; tomarla exigía la alianza interna que produjo el partido. A fines
de 1960 ya no había camino de vuelta para nadie, y la pregunta sobre qué tenía
Castro en la cabeza en enero de 1959 había dejado de importar.

Lo que sí importa para la Parte III es la lección de segundo orden: el problema
de soberanía de Cuba nunca se resolvió, solo se replanteó. El país pasó de
depender de una gran potencia a depender de otra, y en 1990 descubrió lo que eso
había costado. Un diseño para 2026 que no piense a fondo la estructura de la
dependencia externa ---quién compra, quién presta, quién suministra el
combustible--- reproducirá la misma vulnerabilidad con otras banderas.

\section{1959: el año de los dos gobiernos}

Formalmente, Cuba tenía en 1959 un gabinete, un presidente y una constitución de
1940 restaurada y enmendada por una Ley Fundamental. En realidad tenía el
gabinete y tenía el Ejército Rebelde, y donde discrepaban ganaba el Ejército
Rebelde. Castro asumió el cargo de primer ministro en febrero. Urrutia fue
forzado a salir en julio tras criticar la influencia comunista: por televisión,
por Castro, en un acto de teatro político que estableció el método de la década
siguiente.

Dos cosas ocurrieron en 1959 que fijaron todo lo posterior.

\subsection{Los tribunales}

La policía y el ejército de Batista habían torturado y matado, en gran medida con
impunidad, y en enero de 1959 había una demanda popular abrumadora de castigo y
ningún sistema judicial en funciones con legitimidad para satisfacerla. El
gobierno revolucionario juzgó a varios cientos de personas ante tribunales
militares sumarios y las ejecutó, la mayoría en los primeros cuatro meses. La
cifra está en disputa: se citan habitualmente números en torno a 550 para 1959, y
las compilaciones del proyecto Cuba Archive arrojan cifras bastante mayores para
la década completa, mientras que el gobierno nunca ha publicado un conteo.\unv{}
Los juicios fueron públicos, se celebraron ante multitudes y, en muchos casos
individuales, contra personas que demostrablemente habían cometido los actos que
se les imputaban. Fueron también sumarios, no apelaban ante ninguna autoridad
independiente y, en el caso de los aviadores de Batista absueltos, se repitieron
hasta obtener condena por insistencia pública de Castro ---lo que destruyó, en
las primeras semanas, el principio de que un veredicto es definitivo---.

Ambos hechos son ciertos y hay que sostenerlos juntos. Un gobierno revolucionario
que enfrenta una demanda popular genuina de justicia contra un Estado que ha
torturado está ante un dilema real, y no hacer nada no era una opción disponible.
Lo que los tribunales establecieron, sin embargo, fue que en la nueva Cuba la ley
era un instrumento del ejecutivo y no un límite sobre él, y ese precedente jamás
se revirtió. El juicio a Huber Matos, en octubre, usó la misma maquinaria contra
un comandante revolucionario cuyo delito era la discrepancia política, diez meses
después.

\subsection{La reforma agraria}

La Primera Ley de Reforma Agraria del 17 de mayo de 1959 fijó un tope de 30
caballerías por propietario, unas 402 hectáreas, con excepciones de hasta unas
1.340 hectáreas para tierras de alta productividad; expropió el excedente con
compensación en bonos a veinte años al 4,5 por ciento, valorados según las
propias declaraciones fiscales de los dueños; prohibió la propiedad extranjera de
la tierra; y distribuyó tierra a arrendatarios y precaristas en parcelas
inalienables de unas 27 hectáreas. Creó el INRA, el Instituto Nacional de Reforma
Agraria, que en dos años era un gobierno paralelo que manejaba el crédito, la
distribución, la vivienda y buena parte de la economía.

La reforma no era, para los estándares latinoamericanos de la época, extrema
---el tope era generoso, se ofrecía compensación y el propio Estados Unidos
estaría instando a la región a reformas comparables dos años después bajo la
Alianza para el Progreso---. Lo que la convirtió en ruptura fue la valoración y
la moneda. Los propietarios norteamericanos llevaban décadas declarando sus
tierras por una fracción de su valor de mercado para reducir impuestos; la ley
los tomó por su palabra y ofreció el pago en bonos cubanos que nadie creía que
fueran a honrarse. Washington exigió compensación pronta, adecuada y efectiva en
efectivo. La Habana se negó. Ese desacuerdo, más que ningún discurso, es donde
los dos países se separaron.

\section{1960: la mecánica de la ruptura}

El año transcurre como una secuencia de jugadas en la que cada acción de una
parte volvía casi automática la siguiente de la otra.

Febrero: Anastás Mikoyán lleva una delegación comercial soviética a La Habana y
firma un acuerdo para comprar azúcar cubano y suministrar crudo a crédito. Mayo:
se restablecen las relaciones diplomáticas con Moscú. Junio: las refinerías de
propiedad norteamericana y británica en Cuba ---Esso, Texaco, Shell--- se niegan,
por instrucciones, a refinar crudo soviético; Cuba las interviene. Julio: Estados
Unidos cancela el resto de la cuota azucarera cubana. También en julio: la Unión
Soviética anuncia que comprará el azúcar. Agosto: Cuba nacionaliza los servicios
públicos, la compañía telefónica, los centrales y las refinerías norteamericanas.
Octubre: Estados Unidos impone un embargo sobre todas las exportaciones a Cuba
salvo alimentos y medicinas. También en octubre: Cuba nacionaliza esencialmente
el resto de la gran empresa ---la banca, unas 380 grandes firmas nacionales y,
bajo la Ley de Reforma Urbana, toda la vivienda de alquiler, que se convirtió en
propiedad de los inquilinos que la ocupaban, pagando su renta corriente durante
entre cinco y veinte años---.

Esa última medida merece más atención de la que suele recibir, porque sigue
dándole forma a Cuba. La Ley de Reforma Urbana convirtió de la noche a la mañana
a la mayoría de los hogares cubanos en propietarios de su vivienda y abolió el
mercado de alquiler residencial durante seis décadas. Fue enormemente popular; es
la redistribución más duradera que llevó a cabo la revolución; y produjo un parque
habitacional en el que nadie tenía incentivo ni mecanismo legal para invertir en
mantenimiento, y que lleva derrumbándose desde entonces. La Habana pierde
edificios por colapso todos los años. El capítulo~\ref{ch:capital} tiene que
lidiar con la consecuencia, que es que Cuba tiene propiedad de vivienda casi
universal y un parque habitacional catastróficamente deteriorado, y que los dos
hechos son el mismo hecho.

Enero de 1961: Washington rompe relaciones diplomáticas. Febrero de 1962: el
embargo se extiende prácticamente a todo el comercio, bajo la Ley de Comercio con
el Enemigo, donde permanece.

\section{Comienza el éxodo}

Entre 1959 y 1962 se marcharon unos 250.000 cubanos, de 6,6 millones
\citep{eckstein2003}.\unv{} Eran desproporcionadamente la clase profesional y
empresarial: los expropiados, y después sus médicos, abogados, ingenieros y
gerentes. Alrededor de la mitad de los médicos de Cuba se fue en los primeros
años ---unos tres mil de seis mil---, lo que está en el origen de la política
médica de la revolución, porque la respuesta del Estado fue construir escuelas de
medicina a un ritmo que ningún país comparable intentó
\citep{feinsilver1993}.\unv{}

Dos rasgos de esta primera oleada importan más adelante. Fue legal y en buena
medida ordenada, por avión a Miami, y quienes se iban la entendían como temporal
---la mayoría esperaba regresar en cuestión de meses, que es la razón de que se
dejara tanta propiedad al cuidado de familiares, hecho que volverá inusualmente
complicadas las reclamaciones de restitución del capítulo~\ref{ch:capital}---. Y
comprendió la Operación Peter Pan, en la que unos 14.000 niños no acompañados
fueron enviados en avión a Estados Unidos entre 1960 y 1962 por un programa
organizado por la Iglesia católica, sobre la base del rumor de que el Estado se
proponía asumir la custodia de los menores. El rumor era falso. El programa fue
real, muchas familias quedaron separadas durante años, y el episodio envenenó el
registro emocional de la cuestión cubano-americana de un modo que nunca ha
terminado de despejarse.

\section{Playa Girón y octubre de 1962}

El 17 de abril de 1961, unos 1.400 exiliados cubanos, entrenados, armados y
transportados por la CIA, desembarcaron en Bahía de Cochinos. La operación se
había planeado bajo Eisenhower y la heredó Kennedy, que canceló el segundo ataque
aéreo; la cobertura aérea falló, el alzamiento interno esperado no se produjo
---los servicios de seguridad habían detenido a decenas de miles de presuntos
simpatizantes en los días previos--- y la brigada fue derrotada en setenta y dos
horas. Unos 1.200 fueron capturados y más tarde canjeados por alimentos y
medicinas.

La invasión fue la decisión norteamericana más consecuente del período, y su
consecuencia fue la contraria a su propósito. Convirtió la paranoia de seguridad
del gobierno cubano en un hecho demostrado, justificó la movilización permanente,
desacreditó a la oposición interna por asociación con un desembarco extranjero, y
le dio a Castro la ocasión ---el 16 de abril, en el entierro de las víctimas del
bombardeo preliminar--- de declarar por primera vez el carácter socialista de la
revolución. También volvió razonable, en vez de paranoica, la petición cubana de
garantías militares soviéticas.

Dieciocho meses después, en octubre de 1962, el reconocimiento aéreo
norteamericano encontró misiles balísticos soviéticos de alcance medio en
construcción en Cuba. Los trece días que siguieron son lo más cerca que ha estado
el mundo de una guerra nuclear. La solución ---retirada soviética a cambio del
compromiso norteamericano de no invadir, y una retirada discreta de misiles
norteamericanos de Turquía--- se negoció entre Moscú y Washington por encima de
La Habana y contra la furiosa objeción de Castro, que había querido que los
misiles se quedaran y que, en el punto álgido de la crisis, había instado a
Jruschov a considerar un primer golpe si Estados Unidos invadía.

Dos legados. El entendimiento de no invasión es, en la práctica, lo que ha
mantenido vivo al gobierno cubano durante sesenta años, y lo obtuvo la Unión
Soviética y no Cuba. Y la experiencia de haber sido moneda de cambio le enseñó a
la dirigencia cubana una lección sobre el patrocinio de las grandes potencias
sobre la que después declinó actuar durante otros veintiocho años.

\section{La consolidación interna}

Esos mismos tres años construyeron la maquinaria doméstica.

La campaña de alfabetización de 1961 envió al campo durante ocho meses a unos
100.000 voluntarios, la mayoría adolescentes. Reportó haber enseñado a leer a unas
707.000 personas adultas y haber reducido el analfabetismo de alrededor del 23
por ciento a menos del 4.\unv{} El estándar aplicado era elemental y la cifra se
ha cuestionado en los márgenes, pero la realidad de la campaña no está seriamente
en duda ---es visible en los datos de cohortes posteriores--- y sigue siendo lo
más impresionante que hizo la revolución en su primera década. Fue también, y a
la vez, un instrumento de integración política: los brigadistas llevaban un
manual cuyos ejercicios de lectura trataban sobre la revolución, y la campaña
puso a cien mil adolescentes urbanos bajo disciplina estatal en el campo durante
un año.

En junio de 1961 se nacionalizaron las escuelas privadas, poniendo fin a la
educación católica. Ese mismo año se fundaron los Comités de Defensa de la
Revolución, una red cuadra por cuadra que combinaba las campañas de vacunación,
la donación de sangre y la defensa civil con la vigilancia vecinal ---una sola
organización haciendo ambas cosas, que es exactamente la razón de que funcionara
y exactamente la razón de que se la temiera---. La prensa independiente dejó de
existir entre 1960 y 1961 por una mezcla de expropiación, toma por parte de los
sindicatos de tipógrafos y emigración. En marzo de 1962 se introdujo el
racionamiento, la libreta, concebido como medida temporal de emergencia; sigue
vigente en 2026.

Hacia 1962 las organizaciones revolucionarias se habían fusionado en las ORI, y
su primera crisis ---la remoción del viejo comunista Aníbal Escalante en marzo de
1962 por haber llenado el aparato con su propia gente, denunciada por Castro como
``sectarismo''--- estableció quién iba a dirigir el partido. No fueron los
comunistas.

\begin{ledger}{1959--1962}
\emph{Logrado.} Tierra redistribuida a más de 100.000 familias arrendatarias;
alquileres abolidos y viviendas urbanas transferidas a sus ocupantes;
analfabetismo eliminado en la práctica en un solo año; segregación racial del
espacio público y semipúblico terminada por decreto y hecha cumplir; la mitad
rural del país puesta por primera vez al alcance del Estado; fin de la propiedad
extranjera de las alturas dominantes; y una invasión derrotada.

\emph{Costo.} Ejecuciones sumarias por cientos, sin apelación y sin conteo
publicado; los tribunales subordinados al ejecutivo en cuestión de semanas y
nunca restaurados; la prensa independiente extinguida en dos años; toda
organización política distinta de la gobernante disuelta; un cuarto de millón de
personas, incluida la mitad de los médicos, fuera del país; educación privada
abolida; y el comienzo de un aparato de seguridad cuyo alcance a nivel de cuadra
no tenía precedente en el hemisferio. Racionamiento introducido como medida
temporal hace sesenta y cuatro años.
\end{ledger}
